% Solutions to the problems of Chapter 8 (Open Challenges and Future Directions).
% Included by backmatter/solutions.tex through \InputIfFileExists.
% [October 2026 round] New file (blue).
\begin{sol}{prob:ch08:1}
(a) For each sensor $i$ and each detection associated to a confirmed track, the monitor computes the innovation $\boldsymbol{\nu}_{i,k}$, its covariance $\mathbf{S}_{i,k}$ and the NIS $d^2_{i,k}$ of \eqref{eq:nis}; over a sliding window of $N$ samples (typically $N=20$--$50$, i.e.\ $2$--$5$~s at $10$~Hz) it keeps the windowed mean NIS $\bar d^2_i$, the windowed mean innovation (a bias estimate), the fraction of expected returns actually received (detection rate), and, when two sensors observe the same track, the Mahalanobis distance between the single-sensor estimates. (b) The windowed mean NIS is compared with $\chi^2_{Nn_y,1-\alpha}/N$, where $n_y$ is the measurement dimension; $\alpha$ is chosen from the tolerated false-alarm rate per window (for instance $10^{-3}$) and a $k$-out-of-$m$ persistence rule further reduces false alarms at the cost of detection delay. The detection-rate threshold is set from nominal statistics, for instance fewer than half of the expected returns during one second. (c) A reliability factor $\alpha_i\in[0,1]$ is derived as a decreasing function of the excess NIS, e.g.\ $\alpha_i=\min\{1,\ \chi^2_{Nn_y,1-\alpha}/(N\,\bar d^2_i)\}$, and multiplied by a second factor derived from the detection rate; the factor is used both as an evidence-theoretic discount (mass $1-\alpha_i$ is transferred to the frame of discernment) and as a covariance inflation $\mathbf{R}_i/\alpha_i$ in the tracker, as a coordinated heuristic for the two layers of Fig.~\ref{fig:hybrid-pipeline}, not as a proof of calibration. At $\alpha_i=0$ skip that sensor update rather than divide by zero. Monitor pre-gate residuals as well, since gating truncates the nominal NIS distribution; do not tune the reference covariance on the same window and still claim the original false-alarm rate. (d) A blocked sensor shows a collapse of the detection rate while the NIS of the detections that still arrive is normal: the sensor is declared unavailable, the tracker continues with the remaining modalities, and the planner is informed of the reduced field of view and of the larger covariance (a larger protection level). A miscalibrated sensor shows a normal detection rate but a biased innovation whose sign is consistent and, for a temporal offset, grows with the relative velocity (Problem~\ref{prob:ch08:3}): the sensor is discounted, online calibration is triggered, and the planner receives a degraded-integrity flag with the inflated covariance. If either condition persists beyond a pre-defined time and the remaining sensors cannot meet the required integrity, the policy requests a minimal-risk manoeuvre.
\end{sol}

\begin{sol}{prob:ch08:2}
(a) With edges $0,0.2,\dots,1$ the first bin is empty; bin $[0.2,0.4)$ contains two samples (mean confidence $0.285$, accuracy $0.5$, gap $0.215$); bin $[0.4,0.6)$ three samples ($0.517$, $0$, gap $0.517$); bin $[0.6,0.8)$ three samples ($0.717$, $0.667$, gap $0.050$); bin $[0.8,1]$ four samples ($0.900$, $0.750$, gap $0.150$). Hence $\mathrm{ECE}=(2\cdot0.215+3\cdot0.517+3\cdot0.050+4\cdot0.150)/12=0.2275\approx0.23$. (b) The Brier score is $0.2313$; the accuracy is $6/12=0.50$ and the mean confidence $0.656$. (c) The mean confidence exceeds the accuracy, suggesting aggregate over-confidence and motivating a check of $T>1$ for these samples; the optimum must be computed from the labelled likelihood rather than inferred from the two aggregate means alone. Minimising the negative log-likelihood (script \repopath{ch08\_\allowbreak p2\_\allowbreak ece.py} in \repopath{companion\_\allowbreak repo/\allowbreak exercises/\allowbreak ch08/\allowbreak solutions/\allowbreak }) gives $T^\ast\approx1.82$, reducing the NLL from $0.656$ to $0.623$ and the ECE from $0.228$ to $0.132$. (d) With two to four samples per bin, the per-bin accuracy has a standard error of about $0.25$--$0.35$, and the absolute value in the ECE makes the estimate biased upwards for small samples. There is no universal required sample count: report bin counts and confidence intervals, examine sensitivity to binning, and use held-out data. Adaptive bins or smoothing can reduce some estimation problems but also require validation; the exercise illustrates the mechanics. Fitting and scoring $T$ on the same twelve observations gives an optimistic in-sample assessment, not evidence of improved deployment calibration.
\end{sol}

\begin{sol}{prob:ch08:3}
(a) A measurement taken at $t_k-\delta$ but interpreted at $t_k$ reports the range $r(t_k-\delta)=r(t_k)+v\delta$ for a target closing at speed $v$, so the innovation is $\nu_k=v\delta+e_k$ with $e_k\sim\mathcal{N}(0,\sigma^2)$ and $b=v\delta$: $b=0.30$~m for $\delta=20$~ms and $b=0.75$~m for $\delta=50$~ms. (b) $\mathbb{E}[\nu_k^2/\sigma^2]=(b^2+\sigma^2)/\sigma^2=1+b^2/\sigma^2$, i.e.\ $2.0$ and $7.25$. (c) At $50$~ms the expected NIS exceeds the threshold, while individual measurements are rejected probabilistically; since $\nu_k/\sigma\sim\mathcal{N}(b/\sigma,1)$, the rejection probability is $P(\lvert Z+b/\sigma\rvert>1.96)$, which equals $0.17$ at $20$~ms and $0.70$ at $50$~ms, against $0.05$ without offset. (d) The mean-crossing criterion is $1+b^2/\sigma^2>1.571$, i.e.\ $\lvert b\rvert>0.3\sqrt{0.571}=0.227$~m, giving mean-crossing offsets of $15.1$~ms at $15$~m/s and $113.3$~ms at $2$~m/s. For an actual detection-power criterion, $\sum_{k=1}^{20}\nu_k^2/\sigma^2\sim\chi'^2_{20}(20(v\delta/\sigma)^2)$. Solving $P(\chi'^2_{20}(\lambda)>\chi^2_{20,0.95})=0.95$ gives $\lambda=30.72$, hence offsets $24.8$ and $185.9$~ms. The same nominal false-alarm probability thus supports different detection limits depending on the required power. (e) The same offset is both most harmful and most detectable at motorway speeds; in slow urban traffic a $100$~ms offset is less detectable by this window ($0.2$~m bias), although whether that error is tolerable depends on the application. Since $\partial h/\partial\delta=v$, nonzero relative speed gives sensitivity to the offset. If range or an additional constant range bias is unknown, a constant-speed offset is confounded with that bias; a reference or varying velocity is needed to identify them separately. With sufficient excitation it can be appended to the state vector and estimated (augmented-state or unified spatio-temporal calibration), or removed at the source by hardware time-stamping. The script \repopath{ch08\_\allowbreak p3\_\allowbreak timestamp\_\allowbreak offset.py} in the same folder reproduces these numbers by Monte Carlo simulation.
\end{sol}

\begin{sol}{prob:ch08:4}
Feed three trackers with the camera-only, LiDAR-only, and fused detections of the simulation-trained detector and record, over real recordings, (i) the windowed mean NIS per modality, which should be close to one per degree of freedom if the learned covariance transfers (the ratio is a residual diagnostic affected by prediction uncertainty, association, bias, and temporal dependence); (ii) the distribution of the conflict mass $K$ between the camera and LiDAR branches, which should match its simulated distribution; (iii) temporal consistency, measured by the fraction of detections associated to confirmed tracks and by the flicker rate of track existence; and (iv) the distribution of confidence scores, compared with its simulated counterpart by a two-sample test. Treat a modality as having passed this screen, without certifying transfer, if its mean NIS stays within the $\chi^2$ bounds of the window, the shift of the $K$ distribution is below a pre-set threshold, and persistence is not lower than in simulation by more than a margin; otherwise first rescale the covariance (a candidate scalar adjustment fitted on separate calibration data, only after checking that covariance scale rather than bias or dynamics is responsible) and re-test, and investigate association, dynamics, sensor faults, and domain shift. Select covariance recalibration, model adaptation, or retraining using independent labelled validation; unchanged score distributions do not exclude a need for retraining. A small labelled subset adds what label-free statistics cannot see: false negatives (a missed object produces no innovation), the calibration error of the scores, and a sanity check of the thresholds themselves.
\end{sol}

\begin{sol}{prob:ch08:5}
(a) Object identifier and provenance (sender identity, sensor lineage, and the list or hop count of contributing sources); time stamp in a declared time base with the sender's clock accuracy; reference frame and the sender's own pose with its covariance; the definition of the state vector; the full covariance matrix rather than a scalar confidence; an existence probability or a declared mapping from detector outputs to belief masses; a declared reliability or health status; and the motion model used by the sender for prediction. (b) A default covariance that is too small gives the message an excessive weight in the covariance-intersection combination and makes the fused covariance over-optimistic (inconsistent); one that is too large discards the information. Without provenance the same track received directly and through a relaying vehicle is treated as two independent pieces of evidence, the covariance shrinks by a factor of about two, and the estimate becomes confidently wrong, which is exactly the double counting that covariance intersection prevents when the receiver knows that the inputs may be correlated; without provenance the receiver must apply intersection to everything and loses accuracy. (c) Compare time-aligned tracks using the covariance of their difference, $\mathbf{P}_r+\mathbf{P}_s-\mathbf{P}_{rs}-\mathbf{P}_{rs}^{\mathsf T}$. A nominal NIS law requires this covariance and the relevant Gaussian assumptions; summing marginal covariances is justified by independence, not by missing provenance. With unknown correlation, use a conservative bound or an empirically validated discrepancy score. Map validated health evidence to covariance inflation or belief discounting as an explicit policy, including a skip rule at zero trust. Authenticate messages and corroborate trust with independent sources; consistency alone does not prove sender correctness.
\end{sol}

\begin{sol}{prob:ch08:6}
Parse the free-text label onto the frame of discernment $\Theta$ (the known classes plus an ``other'' element); an unambiguous label with self-reported confidence $c$ gives $m(\{A\})=c$ and $m(\Theta)=1-c$, and an ambiguous description gives mass to the corresponding subset (for instance ``some animal'' to the set of animal classes), which is precisely what evidence theory allows and a probability vector does not. The discount factor $\alpha$ is estimated on a held-out set as the empirical precision of the model's labels as a function of its reported confidence, and the resulting discounted assignment $m^\alpha(A)=\alpha\,m(A)$, $m^\alpha(\Theta)=1-\alpha+\alpha\,m(\Theta)$ must be checked separately. Precision at confidence $c$ is not automatically an additional multiplier: if $m(A)=c$ already represents calibrated correctness, multiplying again by precision generally double-discounts it. One option is to fit the complete mapping from model outputs to masses on held-out data, with the intended belief or decision calibration criterion stated explicitly. A large conflict mass $K$ indicates disagreement and should remain visible, as in the transferable belief model. It can arise from semantic error, sensor faults, frame mismatch, or source dependence; classical monitors can help diagnose it but do not uniquely identify the cause. Insert the source in an advisory semantic layer at Fig.~\ref{fig:hybrid-pipeline}, with an explicit planner policy preventing unvalidated semantic outputs from overriding geometrically supported tracks. This authority restriction must be implemented and validated; NIS consistency alone does not prove a track correct.
\end{sol}
